\subsection{Repair and Fresh Extraction with Shared Evidence}
\label{app:reextraction_diagnostics}

We collect 240 further Gemma responses on the same sixty receipt documents:
repair or fresh extraction, each with or without the fixed evidence index.
All arms share source lines, field definitions, the required document ID,
and the completion budget. Repair additionally receives the existing graph.
The primary contrast is indexed repair minus indexed extraction after filtering.
Every response wraps
its JSON object in a Markdown code fence. The frozen strict parser marks all
240 as failures, giving zero F1 in every arm.

After observing this format issue, we add a separate compatibility analysis
that removes only an enclosing code fence, then applies the original schema
checks. It recovers all 240 responses. Table~\ref{tab:reextraction_recovery}
reports this analysis; raw responses and strict scores remain available.
Indexed repair reaches 83.93\% F1 and indexed extraction 61.05\%, a paired
difference of 22.88 points (95\% CI: 14.84 to 31.35).

\begin{table}[t]
\centering
\caption{Same-response diagnostics on sixty receipts. ``Fence-compatible''
removes an enclosing Markdown fence before the unchanged filter.
``ID-normalized'' also assigns the public document ID before filtering,
to separate head errors from field selection. Both are post-hoc analyses.}
\label{tab:reextraction_recovery}
\small
\begin{tabular}{lrrr}
\toprule
Method & Fence-compatible F1 & ID-normalized F1 & Wrong-ID graphs \\
\midrule
Fresh extraction & 0.6333 & 0.7256 & 18/60 \\
Indexed extraction & 0.6105 & 0.7486 & 28/60 \\
Repair & 0.7917 & 0.7917 & 0/60 \\
Indexed repair & 0.8393 & 0.8393 & 0/60 \\
\bottomrule
\end{tabular}
\end{table}

Head errors account for much of the gap: indexed extraction uses a wrong
document ID in 28 graphs, whereas indexed repair uses the correct ID in all
sixty. Assigning the known ID uniformly across all four arms raises indexed
extraction F1 to 74.86\%. Indexed repair still leads by 9.07 points
(95\% CI: 5.22 to 13.10). This diagnostic holds relation values fixed and
shows that the full 22.88-point gap is not solely a field-repair effect.
The comparison measures the value of an available initial graph on reused
receipts. Its construction cost is outside these repair/extraction requests.

\paragraph{Where the index locates useful evidence.}
A separate analysis of the original follow-up localizes its 240 reference
fields using the fixed index. Matching preserves case and punctuation and
normalizes whitespace. A cross-line value counts as covered only when all
lines of one actual source occurrence are included.
Of the 240 values, 188 occur in indexed or nearby lines, 36 occur only
outside them, and 16 do not match the source. Among the 188 covered values,
Simple returns 163 exactly and Evidence Index returns 174. Both return 26
of the 36 values outside the index and none of the sixteen unmatched values.
The net gain of eleven exact fields thus occurs in the covered group.
Source mismatches are kept separate from index misses; these descriptive
associations complement the randomized-context comparisons.
